% NAACL 2027 draft based on the official *ACL style.
%
% Compile inside the official ACL template bundle:
%   pdflatex naacl2027_draft
%   bibtex naacl2027_draft
%   pdflatex naacl2027_draft
%   pdflatex naacl2027_draft
%
% The current official files are maintained at:
%   https://github.com/acl-org/acl-style-files
%
% Before submission:
%   1. Replace every \todo{...} marker.
%   2. Re-check the NAACL 2027/ARR formatting and responsible-NLP checklist.
%   3. Keep [review] for submission; switch to [final] only for camera-ready.

\documentclass[11pt]{article}

\usepackage[review]{acl}
\usepackage{times}
\usepackage{latexsym}
\usepackage[T1]{fontenc}
\usepackage[utf8]{inputenc}
\usepackage{microtype}
\usepackage{inconsolata}
\usepackage{booktabs}
\usepackage{amsmath}
\usepackage{amssymb}
\usepackage{xspace}
\usepackage{xcolor}
\usepackage{url}

\newcommand{\method}{\textsc{TraceGuard}\xspace}
\newcommand{\todo}[1]{\textcolor{red}{[TODO: #1]}}
\newcommand{\bx}{\mathbf{x}}
\newcommand{\by}{\mathbf{y}}
\newcommand{\bz}{\mathbf{z}}

\title{When Answers Stop Leaking and Start Teaching:\\
Ground-Truth-Guided Traces for Self-Improving Reasoners}

% The review option anonymizes this block. Replace it before camera-ready.
\author{Anonymous NAACL 2027 submission}

\begin{document}
\maketitle

\begin{abstract}
High-quality reasoning traces are a key ingredient for improving reasoning
models, yet producing them often requires stronger teacher models or
human-written rationales.  We show that ground-truth answers can instead serve
as privileged guidance for synthesizing high-value, model-native training
traces.  Our pipeline conditions a reasoner on both a problem and its known
answer, allowing ordinary question--answer data to yield substantially stronger
solution trajectories.  This advantage introduces \emph{privileged leakage}:
generated traces may explicitly mention the supplied answer or rely on it as
evidence.  To make these traces valid beyond the privileged generation setting,
we calibrate a
sparse attention-based detector that identifies tokens associated with leakage.
When the detector fires during generation, the decoder rolls back the affected
region and rewrites it using a synchronized clean branch that never observes the
answer.  The resulting ground-truth-guided traces preserve the central benefits
of privileged generation: the known outcome steers generation toward correct
solutions and supports self-verification throughout the reasoning process.
Meanwhile, clean-branch repair removes explicit reliance on privileged context
while preserving fluent and natural reasoning trajectories.  Across
mathematical, scientific, logical, and multi-hop reasoning domains, fine-tuning
on only 1K--2K generated samples consistently improves reasoning performance.
On the MATH domain, the resulting data improves
DeepSeek-R1-Distill-Qwen-14B by up to 14 points.  These results establish
ground-truth-guided trace synthesis as an effective source of reasoning
supervision for bootstrapping self-improving reasoners without external
rationale supervision.
\end{abstract}

\section{Introduction}

Chain-of-thought (CoT) supervision can transfer reasoning behavior to smaller
language models and improve performance on multi-step tasks
\citep{wei2022chain,ho2023reasoningteacher}.  However, creating such supervision
at scale remains difficult.  Unconditioned sampling often produces incorrect
answers, while filtering for correctness discards potentially useful examples
and spends compute on failed trajectories.  A natural alternative is to expose
the known answer while asking the generator to construct a derivation.  This
form of answer-conditioned rationalization is used explicitly in STaR
\citep{zelikman2022star} and is closely related to contrastive answer-conditioned
decoding in SCOTT \citep{wang2023scott}.

Answer conditioning creates a train--test asymmetry.  The generator may mention
that the answer was supplied, verify its work against the reference, or treat
the target as a premise.  We call these behaviors \emph{privileged leakage}.
They are distinct from legitimately deriving and stating the same final answer:
the issue is not the appearance of the answer value itself, but reliance on
information unavailable to the eventual student.  A trace such as ``since the
given answer is 72'' is therefore invalid supervision even when its arithmetic
is correct.  Lexical filters can catch fixed phrases, but leakage can be
paraphrased, distributed across several tokens, or expressed without words such
as \emph{answer} or \emph{ground truth}.

This paper asks whether the generator's own attention can reveal when
privileged information is about to surface.  Our starting observation is that
only a small number of attention heads need to align strongly with human- or
judge-annotated leakage spans.  This agrees with broader evidence that sparse
attention heads can expose task-relevant model behavior
\citep{hung2024attentiontracker}.  We use the mass assigned by these heads to
the privileged prompt span as a token-level signal.  Importantly, detection is
only half of the problem: discarding the full trajectory wastes its valid
prefix, while merely suppressing one token can defer the same leakage to the
next step.

We propose \method, a detect--rollback--regenerate procedure for constructing
reasoning supervision.  Generation maintains two views of the same output
prefix: a \emph{privileged branch} conditioned on the question and answer, and
a \emph{clean branch} conditioned on the question alone.  The privileged branch
normally proposes tokens.  If a proposed token attends too strongly to the
answer span, the decoder backtracks to the beginning of the suspected region
and uses the clean branch to rewrite it.  Both branches then consume the chosen
tokens so that they remain synchronized.  Repair ends only after the detector
observes a configurable safe streak, preventing immediate oscillation back into
the same failure mode.

Our contributions are:

\begin{itemize}
    \item We formulate \emph{privileged leakage} as a distinct data-quality
    problem in answer-conditioned reasoning-trace generation.
    \item We develop a lightweight token-level detector by calibrating a sparse
    set of attention heads against annotated leakage spans.
    \item We introduce online local repair that combines bounded backtracking,
    a clean counterfactual branch, synchronized decoding, and KV-cache reuse.
    \item We provide preliminary detector validation across four Qwen3 model
    scales and an evaluation plan for trace quality, downstream utility, and
    computational overhead.
\end{itemize}

\section{Related Work}

\paragraph{Reasoning data generation and distillation.}
CoT prompting elicits intermediate reasoning from sufficiently capable models
\citep{wei2022chain}, while self-consistency improves answer selection by
sampling multiple reasoning paths \citep{wang2023selfconsistency}.  Fine-tune-CoT
uses teacher-generated rationales as supervision \citep{ho2023reasoningteacher},
and STaR bootstraps rationales by retrying failed questions with the correct
answer supplied \citep{zelikman2022star}.  SCOTT conditions rationale generation
on the target answer through contrastive decoding and emphasizes consistency
between rationales and predictions \citep{wang2023scott}.  These approaches
motivate using target information during data construction; our focus is the
complementary problem of ensuring that the resulting text does not reveal this
privileged construction process.

\paragraph{Faithfulness and internal signals.}
Free-text rationales need not faithfully describe the computation supporting a
prediction, and their surface plausibility alone is therefore insufficient
\citep{jacovi2020faithfulness}.  Work on attention-based prompt-injection
detection shows that selected heads can provide useful inference-time signals
without an auxiliary model call \citep{hung2024attentiontracker}.  We similarly
calibrate sparse heads, but target a different phenomenon: token-level transfer
from a known answer in the prompt into a generated rationale.  We additionally
couple detection to a causal intervention---regeneration without the privileged
context---rather than using the signal only for classification.

\section{Method}
\label{sec:method}

\subsection{Problem Formulation}

Let a supervised reasoning dataset contain pairs $(\bx,\by)$, where $\bx$ is a
problem and $\by$ is its known answer.  We wish to produce a reasoning trace
$\bz=(z_1,\ldots,z_T)$ that is useful for supervised fine-tuning under the clean
input $\bx$.  During data construction, however, we allow the generator
$p_\theta$ to observe a privileged context $c(\by)$, such as
``Given the ground-truth answer is $\by$.'':

\begin{equation}
  z_t \sim p_\theta(\cdot \mid \bx, c(\by), z_{<t}).
  \label{eq:privileged-generation}
\end{equation}

The desired trace may derive and state $\by$, but it should not explicitly
refer to $c(\by)$, cite the supplied target as evidence, or express awareness
that privileged information was provided.  We represent these undesirable
positions with a binary token mask $g\in\{0,1\}^T$ and learn a detector
$D_t$ that estimates whether generation around position $t$ depends on the
privileged span in a way associated with leakage.

\subsection{Constructing Leakage Annotations}
\label{sec:annotations}

We first sample responses from the generator with $c(\by)$ appended to the
problem.  A separate judge receives the privileged prompt and response and
returns exact text spans exhibiting one of four behaviors: (i) directly
referencing the provided answer, (ii) checking a result against it, (iii) using
it as a reasoning constraint, or (iv) discussing the need to hide it.  Merely
deriving or stating the final value is labeled clean.  Character spans are
mapped to response-token spans through tokenizer offsets.  We split annotated
examples into head-calibration, validation, and held-out test sets.

Because judge-generated span boundaries may be imperfect, our primary detector
metric is \emph{span-start recall}: a gold span is detected when its first token
is flagged.  We additionally report span precision, token precision, and
\emph{full recall}, the fraction of examples in which every annotated span is
detected.  \todo{Report annotation protocol, judge model, number of examples,
inter-annotator agreement or a human audit, and the exact data split.}

\subsection{Sparse Attention-Based Leakage Detector}
\label{sec:detector}

Let $A^{\ell,h}_{t,j}$ be the causal attention weight from response token $t$ to
prompt token $j$ at layer $\ell$ and head $h$.  If $I_c$ is the set of token
indices occupied by $c(\by)$, the privileged attention mass is

\begin{equation}
  a^{\ell,h}_t = \sum_{j\in I_c} A^{\ell,h}_{t,j}.
  \label{eq:attention-mass}
\end{equation}

We score each head on the calibration split in two stages.  First, we remove its
within-example background level using a positive median normalization,

\begin{equation}
  \widetilde{a}^{\ell,h}_t =
  \max\!\left(0,
  a^{\ell,h}_t-\operatorname{median}_{u}a^{\ell,h}_u\right).
\end{equation}

Second, we reward both alignment with the gold leakage mask and separation
between positive and negative tokens:

\begin{align}
q_{\ell,h} ={}&
  \cos\!\left(\widetilde{\mathbf{a}}^{\ell,h},\mathbf{g}\right)
  \cdot
  \Bigg[
  \underset{t:g_t=1}{\operatorname{mean}} a^{\ell,h}_t
  -\underset{t:g_t=0}{\operatorname{mean}} a^{\ell,h}_t
  \Bigg]_+ .
\label{eq:head-score}
\end{align}

Scores are averaged over calibration examples, and the top $K$ heads are kept.
Their normalized signals are combined using non-negative calibration scores:

\begin{equation}
  D_t = \frac{\sum_{(\ell,h)\in\mathcal{H}}
  q_{\ell,h}\widetilde{a}^{\ell,h}_t}
  {\sum_{(\ell,h)\in\mathcal{H}}q_{\ell,h}}.
  \label{eq:aggregate}
\end{equation}

We apply a centered rolling maximum of width $w$ to tolerate boundary noise and
flag positions whose score exceeds $\tau$.  The tuple $(K,w,\tau)$ is selected
on the validation split.  Our default selection favors configurations with high
example-level full recall, then balances span recall and precision while
preferring fewer heads and shorter windows.  This produces a model-specific
detector with no learned classifier parameters.

\paragraph{Online score.}
At decode time, only the selected detector layers return attention matrices;
other layers may use a fused attention backend.  This avoids materializing
full-model attention for every token.  The current implementation uses the
calibrated heads' weighted privileged attention mass for the candidate token and
uses $w$ to determine the repair radius.

% IMPORTANT IMPLEMENTATION NOTE (not printed in the paper):
% The current calibration code applies positive-median normalization and a
% centered rolling maximum, whereas the online generator directly thresholds
% aggregated raw attention and uses window_size for backtracking. Before final
% experiments, either align the online computation with Eqs. (3--5), or
% recalibrate tau for the exact online score and update this subsection.

\subsection{Dual-Branch Online Repair}
\label{sec:repair}

For each example, we maintain two autoregressive branches sharing model
parameters but using different prompt caches:

\begin{align}
p_{\mathrm{priv}}(z_t) &=
  p_\theta(z_t\mid \bx,c(\by),z_{<t}),\\
p_{\mathrm{clean}}(z_t) &=
  p_\theta(z_t\mid \bx,z_{<t}).
\end{align}

Both branches always consume the accepted output token, so they condition on an
identical generated prefix even though their prompt contexts differ.  Normal
decoding samples a candidate $\hat z_t$ from $p_{\mathrm{priv}}$ and computes its
detector score.  If the score is below $\tau$, the token is accepted on both
branches.  Otherwise, the method initiates repair.

Let $r=\lfloor(w-1)/2\rfloor$.  Upon detecting leakage at position $t$, we roll
the privileged branch back to $s=\max(1,t-r)$ and discard $z_{s:t}$ together
with the rejected candidate.  The clean branch is restored to the same accepted
prefix.  Tokens covering the backtracked region are sampled from
$p_{\mathrm{clean}}$ and consumed by both branches.  Thereafter, we continue
sampling clean replacements while a privileged candidate remains unsafe.  To
avoid ending repair at an isolated low-score token, generation returns to normal
only after $r$ consecutive privileged candidates pass the detector.  An
alternative comparison uses the Jensen--Shannon divergence between privileged
and clean next-token distributions:

\begin{align}
 m &= \tfrac12(p_{\mathrm{priv}}+p_{\mathrm{clean}}),\\
 \operatorname{JSD}(p_{\mathrm{priv}},p_{\mathrm{clean}})
 &= \tfrac12\operatorname{KL}(p_{\mathrm{priv}}\Vert m) \notag\\
 &\quad +\tfrac12\operatorname{KL}(p_{\mathrm{clean}}\Vert m).
\end{align}

This alternative is useful in ablations but doubles next-token distribution
comparisons and lacks the sparse attribution supplied by the attention detector.

\paragraph{Termination and anti-oscillation.}
Repair is capped at a fixed number of steps.  If repeated detections would
backtrack to the same prefix more than a configured number of cycles, we
suppress further rollback at that location and replace only the newly proposed
token.  These guards make the procedure terminate or fail explicitly instead
of silently looping.  All repair events---detection score, removed tokens,
replacement tokens, and branch choices---are logged for audit.

\subsection{Efficient Decoding}

Naively restarting after every repair would repeatedly prefill the full prompt
and accepted trace.  We instead crop each branch's dynamic KV cache to the
common prefix and recompute only the missing suffix.  During ordinary decoding,
samples with different cache lengths are padded and batched; after a forward
pass, caches are unpadded and returned to per-sample state.  On multi-GPU nodes,
we run one model replica per GPU and distribute examples dynamically.  These
optimizations do not change the detector or token-selection rules, although
different fused kernels and batch shapes may introduce small numerical changes
near a sampling or detection boundary.

\subsection{Exporting Traces for Supervised Fine-Tuning}

The final training example contains the clean user message $\bx$ and the repaired
assistant trace $\bz$; $c(\by)$ is never exported.  We retain the answer value
when it is naturally derived or stated.  Truncated generations may optionally
be excluded, and generation metadata and repair logs are stored separately from
the SFT file.  Thus, the privileged answer acts only as a data-construction aid,
not as student input.

\section{Preliminary Detector Evaluation}
\label{sec:prelim}

Table~\ref{tab:detector} reports the detector metrics currently stored by the
repository for four Qwen3 scales \citep{yang2025qwen3}.  Results are from model-specific detectors and
held-out leakage-annotated responses.  These operating points intentionally
favor recall because a missed span is exported into training data, whereas a
false positive triggers a local clean rewrite.  The high recall therefore comes
with modest span precision, especially for 14B and 32B.

\begin{table}[t]
\centering
\small
\begin{tabular}{lrrrr}
\toprule
Model & Heads & Span P & Span R & Full R \\
\midrule
Qwen3-4B  & 1 & 21.5 & 100.0 & 100.0 \\
Qwen3-8B  & 1 & 38.5 & 99.3  & 98.2 \\
Qwen3-14B & 4 & 8.2  & 99.2  & 98.0 \\
Qwen3-32B & 1 & 13.4 & 99.6  & 98.4 \\
\bottomrule
\end{tabular}
\caption{Preliminary held-out leakage detection results (\%). ``Full R'' is the
percentage of examples for which all annotated leakage spans are detected.
These values validate the detector only; they do not yet establish that online
repair improves final traces or downstream models.}
\label{tab:detector}
\end{table}

\section{Experimental Plan}
\label{sec:experiments}

\todo{Replace this section with completed experiments.}  We propose evaluating
the full pipeline on mathematical, scientific, logical, and multi-hop reasoning
benchmarks.  For each base generator, compare:

\begin{enumerate}
    \item clean generation without the answer;
    \item privileged generation without filtering;
    \item privileged generation with lexical/rule-based filtering;
    \item rejection sampling that discards any detected trace;
    \item \method with attention-based local repair; and
    \item \method with JSD-based repair.
\end{enumerate}

The primary trace-level outcomes should be answer accuracy, leakage rate under a
blinded human audit and an independent judge, retention rate, trace length, and
generation tokens or wall-clock time per accepted example.  To test usefulness
rather than surface cleanliness alone, fine-tune identical student models on
equal-sized trace sets and report clean-input benchmark accuracy.  Detector
ablations should vary the number of heads, calibration size, window length,
threshold, backtracking, safe-streak rule, and clean-branch synchronization.
Cross-model and cross-domain transfer will reveal whether calibration must be
repeated for each generator and task family.

\section{Limitations}

The current evidence validates a detector on automatically annotated leakage
spans but does not yet demonstrate improved downstream reasoning.  The judge may
miss subtle semantic dependence on the supplied answer, and attention is not a
complete causal explanation of generation.  Model-specific heads and thresholds
may not transfer across architectures, checkpoints, domains, prompt templates,
or precision settings.  Recall-oriented thresholds can cause unnecessary clean
rewrites, reducing the advantage of privileged conditioning.  Maintaining two
prompt caches and selectively materializing attention increases memory and
latency relative to ordinary decoding.  Finally, a repaired trace can remain
factually incorrect or unfaithful even when it contains no explicit privileged
leakage.  We therefore treat leakage repair as one data-quality filter, not a
guarantee of correct or faithful reasoning.

\section{Ethical Considerations}

The method is intended to reduce misleading supervision in synthetic reasoning
data.  Nevertheless, automatically generated rationales can convey false
confidence, stereotypes, or unsafe content unrelated to answer leakage.
Releasing generated traces should include their provenance, generator version,
filter settings, and known failure modes.  Any human annotation study should use
appropriate consent, compensation, and data-handling procedures.  \todo{Add the
dataset licenses, compute budget, annotator details, and responsible-NLP
checklist responses used in the final study.}

\section{Conclusion}

We presented \method, a draft framework for generating answer-conditioned
reasoning traces without exporting explicit reliance on the answer.  A sparse
attention detector identifies likely leakage online, and a synchronized clean
branch locally rewrites the affected region while preserving the valid prefix.
Preliminary experiments support the feasibility of high-recall detection across
several model scales.  The decisive next step is an end-to-end evaluation of
whether repaired traces preserve the accuracy benefit of privileged generation
while improving data cleanliness, student performance, and generation
efficiency.

\bibliography{references}

\appendix

\section{Draft Reproducibility Checklist}

The final version should report: model identifiers and revisions; tokenizer and
chat template; attention backend; numerical precision; generation temperature,
top-$p$, token budget, and random seeds; the exact privileged-context template;
detector heads, weights, threshold, and window for every model; repair limits;
annotation prompts; train/validation/test splits; hardware; runtime; and all
filtering decisions.  The implementation already records per-example seeds,
token IDs, repair events, detector-config hashes, and resumable run metadata.

\end{document}
